% ============================================================
% SECCION 4 - RESULTADOS (SINTETIZADA)
% ============================================================

\hypertarget{sec4}{\section{Resultados}}

\begin{keypoint}
\textbf{Hallazgos Principales:}
\begin{itemize}
    \item 72,9\% encontró información falsa en Twitter (plebiscito 2022)
    \item 20.000 bots operando coordinadamente (elección 2021)
    \item Financiamiento corporativo de operaciones de desinformación
\end{itemize}
\end{keypoint}

\sectiondivider

\subsection{Escala del Problema}

\begin{highlight}
La desinformación no es un fenómeno marginal, sino \textbf{operación sistemática, coordinada y bien financiada} que alcanza masivamente a la población.
\end{highlight}

Durante el proceso constituyente (2021-2022):

\begin{enumerate}
    \item \textbf{Fase inicial (jun 2021 - ene 2022)}: Ataques personales contra constituyentes
    \item \textbf{Fase escalada (feb - jun 2022)}: Aumento exponencial sobre contenido del borrador
    \item \textbf{Fase electoral (sep 2022)}: Intensificación en temas sensibles (ahorro, vivienda, salud, educación)
\end{enumerate}

\sectiondivider

\subsection{Temas y Patrones de Ataque}

\begin{description}
    \item[Propiedad:] Falsa afirmación de eliminación de símbolos patrióticos
    \item[Aborto:] Mentira sobre aborto hasta nueve meses
    \item[Vivienda:] Falsedad sobre desalojos estales
    \item[Salud:] Información falsa sobre FONASA y libertad de elección
    \item[Educación:] Aseveración falsa sobre eliminación de colegios particulares
\end{description}

\sectiondivider

\subsection{Ataques Personalizados}

\begin{itemize}
    \item \textbf{Evelyn Matei}: Falsa narrativa sobre Alzheimer, adulteración de videos
    \item \textbf{Elisa Loncón}: Imagen falsa con Augusto Pinochet
    \item \textbf{Gabriel Boric}: Narrativa sobre drogas (exámenes negativos presentados)
    \item \textbf{Mujeres periodistas}: \textbf{76\% de ataques bot} dirigidos a mujeres (Mónica Rincón, Monserrat Álvarez, Diana Levín)
\end{itemize}

\sectiondivider

\subsection{Financiamiento Identificado}

\begin{center}
\begin{tabular}{l>{\centering\arraybackslash}m{3cm}>{\centering\arraybackslash}m{2.5cm}}
\toprule
\cellcolor{headercolor}\textcolor{white}{\textbf{Entidad}} & \cellcolor{headercolor}\textcolor{white}{\textbf{Controlador}} & \cellcolor{headercolor}\textcolor{white}{\textbf{Actividad}} \\
\midrule
\rowcolor{rowcolor}Plan Vital & Grupo Generali (Italia) & Propaganda \\
Provida & Grupo MetLife (NY) & Bots \\
\rowcolor{rowcolor}AFP Habitat & Grupo Prudencial (NJ) & Influencers \\
Currum & Grupo Principal (Iowa) & Deepfakes \\
\bottomrule
\end{tabular}
\end{center}

\sectiondivider

\subsection{Impacto en Percepciones}

\textbf{Polarización Subjetiva:}
\begin{itemize}
    \item Izquierda cree 73\% derecha quiere aborto prohibido; realidad: 40\% (brecha: 33 pts)
    \item 56\% sentía ciudadanía poco/nada informada en plebiscito
    \item 25\% de mensajes WhatsApp político contiene desinformación
\end{itemize}

\sectiondivider

\subsection{Patrones por Colectivo Político}

Porcentaje de afirmaciones verificadas como falsas/imprecisas:

\begin{itemize}
    \item \textbf{Derecha:} 32\% (mayor concentración en falsedades)
    \item \textbf{Independientes:} 32\% (mayor proporción de imprecisiones)
    \item \textbf{Izquierda:} Menor concentración
    \item \textbf{Centro:} Menor concentración
\end{itemize}

\pagedivider
